%Neural Machine Translation (NMT) is a simple new architecture for translating
%texts from one language into another
%\cite{sutskever2014sequence,cho14}. NMT is a single deep 
%neural network that is trained end-to-end, holding several advantages such as the
%ability to capture long-range dependencies in
%sentences, and generalization to unseen texts. Despite being relatively new, NMT has already
%achieved state-of-the-art translation results for several language pairs 
%including English-French \cite{luong2015addressing}, English-German
%\cite{jean2015using,luong2015effective,luong15iwslt,sennrich16mono}, English-Turkish
%\cite{sennrich16mono}, and English-Czech
%\cite{jean15wmt,luong16char}. 
%Figure~\ref{fig:nmt_simple} gives an example of an NMT system.

%\begin{figure}
%\centering
%\includegraphics[width=0.4\textwidth]{img/6-2_nmt_simple} %, trim = 78mm 80mm 133mm 70mm,clip
%\caption{A simplified diagram of NMT.}
%\label{fig:nmt_simple}
%\end{figure}


While NMT has a significantly smaller memory footprint than traditional phrase-based approaches (which need to store gigantic phrase-tables and
language models), the model size of NMT is still prohibitively large for mobile devices.
For example, a recent state-of-the-art NMT system requires over 200 million
parameters, resulting in a storage size of hundreds of megabytes
\cite{luong15attn}. 
Though the trend for bigger and deeper neural networks has brought great progress, it has also introduced over-parameterization, resulting in long running times, overfitting, and the storage size issue discussed above. 
A solution to the over-parameterization problem could potentially aid all three issues, though the first (long running times) is outside the scope of this paper.

\biformat{Our contribution.}
In this section, we investigate the efficacy of weight pruning for NMT as a means of compression.
We show that despite its simplicity, magnitude-based pruning with retraining is highly effective, and we compare three magnitude-based pruning schemes --- \textit{class-blind}, \textit{class-uniform} and \textit{class-distribution}.
Though recent work has chosen to use the latter two, we find the first and simplest scheme --- \textit{class-blind} --- the most successful.
We are able to prune 40\% of the weights of a state-of-the-art NMT system with negligible performance loss, and by adding a retraining phase after pruning, we can prune 80\% with no performance loss.
Our pruning experiments also reveal some patterns in the distribution of
redundancy in NMT. In particular, we find that higher layers, attention and softmax weights are the most important, while lower layers and the embedding weights hold a lot of redundancy. 
For the Long Short-Term Memory (LSTM) architecture, we find that at lower layers the parameters for the input are most crucial, but at higher layers the parameters for the gates also become important.
